\bookchapter{Bab 4: Pasang dan Urus Aplikasi — “Pakej Kemahiran” Telefon}

\chapterintro{Aplikasi menjadikan telefon mampu melakukan apa sahaja, tetapi perlu diurus dengan baik. Sama seperti anda menyusun tempat duduk dan memperuntukkan kebenaran untuk pekerja baharu, memasang App juga memerlukan satu set proses piawaian.}

\booksection{Satu, Kedai Rasmi: “Supermarket Aplikasi” Paling Selamat}

\begin{itemize}
  \item \textbf{iOS}: Pasaran Aplikasi Huawei.
  \item \textbf{Android}Prioritaskan muat turun hanya dari kedai rasmi!
  \item \textbf{HarmonyOS}Jenama tempatan utama (Huawei, Xiaomi, OPPO, vivo, Honor) semuanya mempunyai
\end{itemize}

Mod Tulen\textbf{atau}

\booksection{Dua, Mod Tulen / Perlindungan: Aktifkan “Dinding Anti-Virus” Telefon}

terbina dalam, ini adalah suis keselamatan peringkat sistem.\textbf{Sangat disyorkan untuk semua pengguna mengaktifkannya}.\textbf{1. Apakah sebenarnya ia?}Fungsi teras Mod Tulen (juga dipanggil “Mod Perlindungan”) adalah\textbf{menghalang pemasangan App pihak ketiga}:

\booksubsection{Apakah sebenarnya ia?}

, pakej pemasangan yang dimuat turun dari pelayar, kod QR, perisian sembang dan lain-lain saluran akan disekat secara langsung, menghapuskan perisian hasad dari puncanya.\textbf{2. Kumpulan Pengguna Yang Sesuai}Pengguna warga emas\textbf{: Tidak dapat membezakan antara iklan “rasmi” dan “palsu”, Mod Tulen adalah jimat pelindung yang paling selamat.}Pengguna yang tidak peka terhadap kebenaran aplikasi

\booksubsection{Kumpulan Pengguna Yang Sesuai}

\begin{itemize}
  \item \textbf{Semua pengguna yang ingin meningkatkan keselamatan peranti}: Walaupun peminat, lebih baik diaktifkan daripada tidak — jika benar-benar perlu, boleh dimatikan buat sementara.
  \item \textbf{3. Laluan Pengaktifan (Contoh Huawei/Xiaomi)}Huawei/HarmonyOS
  \item \textbf{: [Tetapan] > [Sistem \& Kemas Kini] >}[Mod Tulen]
\end{itemize}

\booksubsection{Laluan Pengaktifan (Contoh Huawei/Xiaomi)}

\begin{itemize}
  \item \textbf{Xiaomi}: [Tetapan] > [Perlindungan Privasi] >\textbf{[Lindungi Privasi]}atau [Tetapan] > [Tetapan Aplikasi] > [Pemasangan Aplikasi], matikan "Pasang Aplikasi Sumber Tidak Diketahui".
  \item \textbf{OnePlus}: [Tetapan] > [Keselamatan] >\textbf{[Pasang Aplikasi Sumber Luaran]}, matikan kebenaran satu persatu.
  \item \textbf{Apple}Kedai aplikasi bukan "hiasan", ia mempunyai talian pengeluaran keselamatan yang lengkap di sebaliknya:\textbf{Imbasan virus}: Setiap App akan diimbas secara gabungan oleh berbilang enjin anti-virus sebelum disenaraikan, mengenalpasti ciri trojan dan virus yang diketahui.
\end{itemize}

\booksubsection{Mekanisme Pengesanan Keselamatan Kedai Aplikasi}

: Kedai akan menyemak sama ada kebenaran yang dipohon oleh App “munasabah” — App lampu suluh meminta kebenaran senarai kenalan? Terus ditolak atau dinyahsenaraikan.

\begin{itemize}
  \item \textbf{Pengesahan pembangun}: Penyenaraian memerlukan kelayakan perniagaan, pengenalan diri atau pengesahan kad kredit. Jika ada masalah, boleh dikesan, bukan "App haram".
  \item \textbf{Pemantauan berterusan}: App yang telah disenaraikan juga akan diimbas semula secara berterusan, tingkah laku hasad yang baru ditemui akan dinyahsenaraikan secara paksa.
  \item \textbf{Tiga, Sumber Pihak Ketiga dan "Parasit Elektronik": Risiko Sangat Tinggi!}Memuat turun App dari halaman web atau pautan yang tidak jelas mungkin mengandungi virus, mencuri privasi, dan kecurian pembayaran. Di sini kami khusus menyebut satu konsep baharu:
  \item \textbf{"Parasit Elektronik"}.
\end{itemize}

\booksection{Tiga, Sumber Pihak Ketiga dan “Parasit Elektronik”: Risiko Sangat Tinggi!}

Ia merujuk kepada\textbf{perisian “versi retak” yang ditanam trojan}— anda sangka “jimat wang”, sebenarnya “kena jebak”. Bentuk yang biasa:

\textbf{“Versi retak”, “versi percuma”, “versi tanpa iklan” perisian berbayar;}\par

Alat kelabu seperti “pemboost like”, “perebut sampul merah”, “penipu permainan”;\textbf{App “versi dalaman” yang dikatakan disyorkan oleh “iklan kecil” dalam sembang kumpulan.}Program-program ini kelihatan boleh digunakan, tetapi di sebaliknya mungkin

\begin{itemize}
  \item mencuri senarai kenalan, membaca SMS, mengintai mikrofon, melombong di latar belakang
  \item . Mod Tulen adalah
  \item garis pertahanan pertama
\end{itemize}

terhadap mereka — terus menghalang anda daripada memasangnya.\textbf{Cadangan garis pertahanan kedua:}Imbasan virus penuh berkala\textbf{: Gunakan pusat keselamatan terbina dalam telefon (seperti Pengurus Telefon Huawei, Pusat Keselamatan Xiaomi) untuk mengimbas sekali sebulan.}Berhati-hati memberikan kebenaran sensitif

\booksubsection{Cadangan garis pertahanan kedua:}

\begin{itemize}
  \item \textbf{Perhatikan penggunaan kuasa dan data yang tidak normal}: Telefon tiba-tiba panas, bateri cepat habis, data berjalan pada waktu malam — kemungkinan besar terdapat "parasit" tersembunyi di latar belakang.
  \item \textbf{Jika terpaksa memasang App pihak ketiga:}Pastikan sumber benar-benar boleh dipercayai.
  \item \textbf{Baca dengan teliti keperluan kebenaran sebelum memasang.}Imbas dengan App keselamatan selepas pemasangan.
\end{itemize}

Cadangan muktamad: Dalam tetapan sistem, matikan kebenaran “Pasang Aplikasi Sumber Tidak Diketahui”.

\begin{enumerate}
  \item Empat, Pasang, Kemas Kini, Nyahpasang
  \item Pemasangan
  \item : Klik "Dapatkan" atau "Pasang" di kedai, selepas dipasang, buka pertama kali akan meminta kebenaran.
\end{enumerate}

Kemas kini

\booksection{Empat, Pasang, Kemas Kini, Nyahpasang}

\begin{itemize}
  \item \textbf{Mesti kemas kini!}Utamanya untuk membaiki kelemahan keselamatan. Disyorkan mengaktifkan kemas kini automatik.
  \item \textbf{Nyahpasang}: Tekan lama ikon dan seret ke kawasan nyahpasang.\textbf{Lima, Pengurusan Kebenaran: Anda Adalah "Bos" App}meminta kebenaran (seperti senarai kenalan, lokasi) kadang-kadang munasabah (peta memerlukan lokasi), kadang-kadang tidak munasabah (lampu suluh memerlukan senarai kenalan).
  \item \textbf{Pilihan kebenaran}: "Sentiasa benarkan", "Benarkan hanya semasa digunakan", "Larang", "Tanya setiap kali".
\end{itemize}

\booksection{Lima, Pengurusan Kebenaran: Anda Adalah “Bos” App}

:

\begin{itemize}
  \item \textbf{Minimakan: Jangan berikan jika tidak perlu.}Audit selepas itu: Semak secara berkala siapa yang "melampaui kebenaran" dalam tetapan kebenaran.
  \item \textbf{Prinsip pengurusan}:
  \begin{enumerate}
    \item Senarai kenalan/SMS: Boleh digunakan untuk menghantar SMS spam.
    \item Lokasi masa nyata: Boleh dijejak pergerakan anda.
    \item Kamera/Mikrofon: Boleh merakam secara senyap.
  \end{enumerate}
\end{itemize}

\begin{tcbwarn}{Akses fail: Boleh membaca fail peribadi anda.}
\begin{itemize}
  \item Enam, Unsur Pembentukan Pakej Pemasangan Perisian
  \item Pernahkah anda terfikir, apakah kandungan sebenar pakej pemasangan yang dimuat turun dari kedai aplikasi? Memahami struktur pakej pemasangan boleh membantu anda lebih memahami bagaimana App berfungsi, dan juga mengetahui dari mana perlu disiasat apabila menghadapi masalah.
  \item 1. Struktur Fail APK Android
  \item APK (Android Package Kit) pada asasnya adalah arkib ZIP termampat, selepas dinyahmampat, anda boleh melihat direktori dan fail teras berikut:
\end{itemize}

\end{tcbwarn}

\booksection{AndroidManifest.xml}

: “Kad pengenalan” dan “manual arahan” aplikasi. Mencatatkan nama pakej aplikasi, nombor versi, kebenaran yang diperlukan, versi sistem yang disokong, Activity masuk dan lain-lain maklumat teras. Sistem membaca fail ini terlebih dahulu semasa memasang aplikasi.

\booksubsection{Struktur Fail APK Android}

: Kod teras aplikasi. Kod sumber Java/Kotlin selepas dikompil menjadi baitkod Dalvik, disimpan dalam fail ini. Satu aplikasi mungkin mempunyai berbilang fail dex (seperti classes2.dex, classes3.dex).

\begin{itemize}
  \item \textbf{resources.arsc}: Jadual indeks sumber yang telah dikompil. Mencatatkan ID semua sumber dan hubungan yang sepadan.
  \item \textbf{Direktori res/}: Folder sumber. Mengandungi imej, fail susun atur, rentetan, definisi warna dan lain-lain. Imej resolusi berbeza diletakkan dalam folder drawable yang berbeza (seperti drawable-hdpi, drawable-xhdpi).
  \item \textbf{Direktori assets/}: Folder sumber asal. Menyimpan sumber yang tidak dikompil, seperti fail halaman web, fon, audio dan lain-lain. App boleh membacanya melalui AssetManager semasa berjalan.
  \item \textbf{Direktori lib/}: Folder pustaka tempatan. Menyimpan fail .so yang telah dikompil (hasil kompilasi kod C/C++), dikelaskan mengikut seni bina CPU (seperti arm64-v8a, armeabi-v7a). Permainan dan App yang memerlukan pengiraan berprestasi tinggi biasanya mempunyai direktori ini.
  \item \textbf{Direktori META-INF/}: Folder maklumat tandatangan. Mengandungi tiga fail MANIFEST.MF, CERT.SF, CERT.RSA, digunakan untuk mengesahkan integriti dan tandatangan APK.
  \item \textbf{2. Struktur Fail IPA iOS}IPA (iOS App Store Package) juga merupakan arkib termampat, strukturnya sedikit berbeza:
  \item \textbf{Direktori Payload/}: Lokasi kandungan teras. Terdapat folder .app di dalamnya, mengandungi semua kandungan aplikasi.
\end{itemize}

\booksubsection{Struktur Fail IPA iOS}

: Fail konfigurasi aplikasi, serupa dengan AndroidManifest.xml Android. Mencatatkan ID aplikasi, versi, keperluan kebenaran, peranti yang disokong dan lain-lain.

\begin{itemize}
  \item \textbf{Fail boleh laksana}: Program utama dengan nama yang sama dengan folder .app, adalah fail binari format Mach-O yang telah dikompil.
  \item \textbf{Direktori Frameworks/}: Pustaka dinamik dan rangka kerja yang dikehendaki oleh aplikasi.
  \item \textbf{Fail sumber}: Imej, audio, video dan lain-lain sumber, biasanya disimpan dalam folder .app.
  \item \textbf{Direktori \_CodeSignature/}: Maklumat tandatangan kod, digunakan untuk mengesahkan integriti dan sumber aplikasi.
  \item \textbf{3. Mekanisme Pengesahan Tandatangan}Mengapa App yang dimuat turun dari kedai rasmi lebih selamat? Pengesahan tandatangan adalah kuncinya.
  \item \textbf{Tandatangan Android}: Pembangun menandatangani APK dengan kunci peribadi, sistem mengesahkan dengan kunci awam semasa pemasangan. Jika APK diubah suai (seperti menanam virus), tandatangan akan menjadi tidak sah, sistem akan menolak pemasangan. Google Play juga memerlukan App Signing, Google menyimpan kunci tandatangan, melindungi keselamatan aplikasi dengan lebih lanjut.
\end{itemize}

\booksubsection{Mekanisme Pengesahan Tandatangan}

: Semakan Apple lebih ketat. Setiap IPA mesti ditandatangani dengan sijil yang dikeluarkan oleh Apple, jika tidak, ia tidak boleh dipasang. Inilah sebabnya iOS hanya boleh memasang aplikasi dari App Store atau sijil perniagaan. Aplikasi yang ditandatangani sijil perniagaan mempunyai tempoh sah yang terhad, dan Apple boleh membatalkannya bila-bila masa.

\begin{itemize}
  \item \textbf{Fungsi tandatangan}:
  \item \textbf{Mengesahkan sumber: Mengesahkan aplikasi benar-benar dari pembangun yang didakwa.}Mengesahkan integriti: Memastikan aplikasi tidak diubah suai semasa penghantaran.
  \item \textbf{Mencegah penyamaran: Aplikasi hasad tidak boleh menyamar sebagai aplikasi lain.}
  \begin{enumerate}
    \item Anda mungkin mendapati saiz App yang dipaparkan di kedai aplikasi berbeza dengan ruang yang digunakan selepas pemasangan. Ini kerana:
    \item Saiz muat turun
    \item : Saiz APK/IPA selepas pemampatan.
  \end{enumerate}
\end{itemize}

\booksubsection{Saiz Pakej Pemasangan dan Pengoptimuman}

: Saiz selepas nyahmampat, biasanya jauh lebih besar daripada saiz muat turun.

\begin{itemize}
  \item \textbf{Penggunaan masa jalan}: App juga menghasilkan cache dan fail data semasa berjalan, menggunakan lebih banyak ruang.
  \item \textbf{Kaedah pengoptimuman App moden:}App Bundle (Android)
  \item \textbf{: Teknologi pengedaran dinamik Google Play. Pengguna hanya memuat turun sumber yang sesuai untuk peranti mereka semasa memuat turun (seperti hanya memuat turun imej resolusi skrin yang sepadan), mengurangkan saiz muat turun dengan ketara.}App Thinning (iOS)
\end{itemize}

: Teknologi yang serupa, Apple secara automatik memotong sumber yang tidak diperlukan mengikut peranti.

\begin{itemize}
  \item \textbf{Muat turun atas permintaan}: Aplikasi besar seperti permainan boleh memuat turun pakej teras terlebih dahulu, kemudian memuat turun pakej sumber kemudian.
  \item \textbf{Pengetahuan Kecil: Mengapa sesetengah App mempunyai pakej pemasangan yang sangat besar?}App jenis permainan biasanya paling besar kerana mengandungi banyak tekstur definisi tinggi, model 3D, dan fail audio. App sosial juga agak besar kerana perlu menyokong pelbagai modul fungsi. App utiliti biasanya paling kecil. Jika anda mendapati App utiliti tertentu sangat besar, berhati-hati sama ada ia mengikat fungsi yang tidak perlu.
  \item \textbf{Sesi Amali: Semak Kebenaran 3 App}Pilih 3 App yang kerap digunakan (seperti sosial, membeli-belah, utiliti), masuk ke halaman pengurusan kebenaran aplikasi dalam [Tetapan], semak satu persatu, matikan kebenaran sensitif yang tidak perlu atau tukar kepada "Benarkan hanya semasa digunakan".
\end{itemize}

\begin{tcbinfo}{Bab 5: Buat Telefon Lebih Memahami Anda — Tetapan Asas}
Luangkan beberapa minit untuk tetapan, peningkatan pengalaman sangat besar. Tetapan ini seperti “menala” telefon, menjadikannya lebih sesuai dengan tabiat penggunaan anda.

\end{tcbinfo}

\begin{tcbpractice}{Satu, Bunyi dan Getaran}
Tetapkan nada dering dan bunyi pemberitahuan yang anda suka. Gunakan mod senyap (mesyuarat) dan mod jangan ganggu (tidur) dengan fleksibel, dunia serta-merta tenang.

\end{tcbpractice}
